% Layout vs. Schematic (LVS) —— 由 work/build.py 自动生成，勿手改；改译文请改 work/translations/
\chapter{Layout vs. Schematic (LVS)}\label{ch:lvs}

\section{LVS 概览}\label{sec:lvs_overview}

\subsection{LVS 脚本的基本用法}

从 0.26 版起，KLayout 把 LVS 作为内置功能支持。LVS 是版图验证中的一个重要步骤：它确保画出来的电路与设计意图中的 schematic 一致。

基本做法很简单：先分析输入的版图，从中提取出 netlist；再把这份 netlist 与参考 netlist（schematic）比对。若两者等价，电路就很可能按预想的方式工作。

除版图之外，LVS 脚本还需要一份 schematic netlist。目前 KLayout 可以读取 SPICE 格式的 netlist，读取器也支持一定程度的配置，因此可望消化相当一部分实用的 SPICE netlist。

基本思路虽简单，细节却相当复杂。本文档旨在说明 KLayout 为实现 LVS 所提供的方案，以及这一过程带来的各种约束。

KLayout 的 LVS 与 DRC 脚本一样集成在 Macro Development IDE 中；事实上 LVS 是 DRC 脚本的扩展，DRC 的全部功能在 LVS 脚本中同样可用。netlist 提取在 DRC 框架内完成，只是该框架被赋予了识别器件与连接关系的能力。

LVS 脚本在宏编辑器 IDE 中创建、编辑与调试，统一在 \ui{LVS} 标签页中管理。IDE 的更多细节见第 原文对应章节 节；DRC 脚本的入门见第 \ref{sec:drc_basic} 节。

LVS 脚本的 XML 形式用 \ui{.lylvs} 后缀（对应 DRC 的 \ui{.lydrc}），纯文本形式用 \ui{.lvs}（同 \ui{.drc}）。与 DRC 脚本一样，LVS 脚本也能在批处理模式下独立执行，参见第 \ref{sec:drc_basic} 节中“把 KLayout 当作独立的 DRC 引擎”一节。

\subsection{KLayout 的 LVS 实现}

KLayout 内部的 LVS 实现在连接关系、器件识别与输入输出通道方面都设计得高度灵活，要点如下：

\begin{itemize}
\item 尽量少做假设：KLayout 力图少设前提——它不要求 label（虽然 label 有帮助），也不要求特定的层次结构、Cell 名或几何形状。netlist 提取完全由版图中的 Polygon 得出，label 与 Cell 层次只是提供额外信息。
\item 层次化分析：KLayout 具备层次化版图处理引擎来支持层次化 LVS。所谓层次化处理，是指布尔运算尽可能在局部 Cell 环境中完成；相应地，器件也在其所属的版图 Cell 内被识别。
\item 层次稳定：KLayout 不会改动版图的层次结构，也不会引入 variant——至少对布尔运算及部分其它运算是如此。这样即便做过层次化 DRC 运算，版图与 schematic 的层次之间仍能保持对应。只有少数分析会引入 variant。
\item 灵活的引擎：netlist 生成引擎在器件识别与连接提取方面高度灵活。首先，几乎所有 DRC 功能都能用来推导用于器件构造与连接提取的中间图层；其次，器件识别可以用脚本定制。
\item 灵活的 I/O：netlist 是 KLayout 的对象树，其组成（net、device、circuit、subcircuit 等）都完整映射为脚本对象（主类见 API 文档）。因此 netlist 可以在 LVS 脚本内或其它场合被分析与操纵。
\item 界面集成：KLayout 提供 netlist 提取结果与 LVS 报告（cross reference、错误）的浏览器。
\end{itemize}

\subsection{术语}

KLayout 使用一套专门的术语，说明如下：

\begin{itemize}
\item Circuit（电路）：由器件、pin 与 subcircuit 连接而成的图。图的节点是连接至少两个元素的 net。若由版图导出，一个 circuit 对应某个特定的版图 Cell。
\item Abstract circuit（抽象电路）：清空了内部结构的 circuit。这类 circuit 没有 net，只定义 pin，本质上是“黑盒”，LVS 需把其内部结构视为无关（don't care）。抽象电路适合用来表示那些不在版图中出现、但存在于 schematic 中的子单元。
\item Pin（引脚）：circuit 与外界连接的点。circuit 可以以 subcircuit 的形式嵌入别的 circuit；连到这些 subcircuit 的 pin 上的 net 会传播进该 subcircuit，并在那里连接更多元素。pin 通常挂在一个 net 上，某些情况下可以挂多个。
\item Subcircuit（子电路）：嵌入另一个 circuit 中的 circuit。同一个 circuit 可以被多次使用，因此多个 subcircuit 可以引用同一个 circuit。若由版图导出，一个 subcircuit 对应某个特定的 Cell Instance。
\item Device（器件）：描述原子功能单元的 n 端实体，可以是有源器件（晶体管），也可以是无源器件（电阻、电容）。
\item Device class（器件类）：某一类型的器件。器件类有“种”（kind）之分，每种之下可以有多个类。例如 MOS 晶体管的种是 \ui{MOS4}（四端 MOS 晶体管），而在 CMOS 工艺中通常至少有 \ui{NMOS} 与 \ui{PMOS} 两个类。
\item Device extraction（器件提取）：检测器件并在导电区域与器件本体之间建立连接的过程，这些连接最终构成器件的 terminal。
\item Device combination（器件合并）：把同类多个器件组合成单个器件的过程。例如串联电阻可以并成一个；更重要的是，并行的 MOS 晶体管（“指状”晶体管）会合并成单个器件。
\item Terminal（端子）：器件的 pin。terminal 通常按功能命名，例如 MOS 晶体管的栅极叫 \ui{G}。
\item Connectivity（连接关系）：对工艺叠层中导电区域的描述。图层之间存在层内与层间连接：层内连接指同层上相互接触的 Polygon 构成一个连通（即导电）区域；层间连接则由 via 之类建立。
\item Netlist（网表）：由 circuit 与 subcircuit 构成的层次结构。netlist 通常有一个顶层 circuit，其它 circuit 通过 subcircuit 从它被调用。
\item Extracted netlist（提取网表）：由版图导出的 netlist。有时该词指带寄生元件（由连线几何导出的电阻、电容）的网表；在 KLayout 的 LVS 语境中，它指纯粹的连通性 netlist。
\item Schematic（原理图）：作为 LVS 参考的 netlist。“schematic”指的是设计者画出的 netlist，它由物理实现过程转成版图；LVS 则把版图转回 extracted netlist 并与 schematic 比对。
\item Annotated layout 与 Net geometry（标注版图与 net 几何）：把属于各 net 的 Polygon 归集起来。circuit 中的每条 net 都由一组 Polygon 表示，对应原有连线几何与器件 terminal 的几何。由于 net 可以经 pin 传播到 subcircuit，net 及其几何标注也会随之传播。
\item Layout to netlist database（L2N DB）：把 extracted netlist 的信息与标注版图合在一起的数据结构。L2N 数据库可用于可视化 net、从已知位置探测 net，以及做其它分析与处理；有相应的 API 供访问。
\item Cross reference（交叉引用）：LVS netlist 比对（pairing）中两个 netlist 之间相互匹配的对象列表，同时也列出不匹配项与“不精确配对”。不精确配对指并不完全一致、但仍很可能应视为同一对象的一对。
\item LVS database（LVS 数据库）：L2N 数据库、schematic netlist 与 cross reference 的组合，是 LVS 结果的完整映像；有相应 API 供访问其中元素。
\item Labels（标签）：画在版图中的 Text 对象，用来在特定图层的特定位置打上文字。label 最典型的用途是给 net 命名：若 label 参与连接关系，由它形成的 net 就会以该文字作为 net 名。
\end{itemize}

\section{LVS 入门}\label{sec:lvs_intro}

\subsection{LVS 入门}

为介绍 LVS 功能，我们取最简单的 CMOS 结构：两管反相器。

\subsubsection{版图}

该反相器由两个 MOS 晶体管构成。单个晶体管由一块有源区（ACTIVE 层上的矩形）和横跨其上的栅极（POLY 层）组成，栅极在源、漏区（栅极左右两侧）之间形成沟道；CONTACT 提供与上层金属的连接。

实际版图做成标准单元的形式，多个标准单元可以沿水平方向排成一行。电源轨做在第二层金属上：上方为 VDD，下方为 VSS。n-well 延伸到单元顶部，并应与相邻单元的 well 区域相连：

\klfig[0.8\textwidth]{inv.png}{两管反相器版图}

\subsubsection{schematic}

可以为这个反相器画一个简化 schematic（左图），也可以画一个更接近实际、包含晶体管体电位（bulk）的 schematic（右图）。让晶体管体区保持在确定电位以避免闩锁效应很重要，因此这些 terminal 也需要各自的 pin。

\klfig[0.8\textwidth]{inv_schematic.png}{反相器 schematic：简化版（左）与含 bulk 连接的版本（右）}

做 LVS 首先需要一份参考 schematic，即与带 bulk 连接的 schematic 相对应的 SPICE netlist：

\begin{lstlisting}
* Simple CMOS inverer circuit (inv.cir) .SUBCKT INVERTER VSS IN OUT NWELL SUBSTRATE VDD Mp VDD IN OUT NWELL PMOS W=1.5U L=0.25U Mn OUT IN VSS SUBSTRATE NMOS W=0.9U L=0.25U .ENDS
\end{lstlisting}

我们要分析的电路是一个被嵌入更大电路中的 Cell，因此把反相器描述成 subcircuit 比较合适。若 netlist 中只有 subcircuit，KLayout 会分析该 circuit；否则会把全局定义视为主电路。

\subsection{LVS 脚本示例}

用于比对上述版图与 schematic 的 LVS 脚本如下（细节见第 \ref{sec:lvs_intro} 节）：

\begin{lstlisting}
# LVS script (demo technology, KLayout manual) # Preamble: deep # Reports generated: report_lvs # LVS report window # Drawing layers: nwell = input(1, 0) active = input(2, 0) pplus = input(3, 0) nplus = input(4, 0) poly = input(5, 0) contact = input(6, 0) metal1 = input(7, 0) metal1_lbl = labels(7, 1) via1 = input(8, 0) metal2 = input(9, 0) metal2_lbl = labels(9, 1) # Bulk layer for terminal provisioning: bulk = polygon_layer # Computed layers: active_in_nwell = active & nwell pactive = active_in_nwell & pplus pgate = pactive & poly psd = pactive - pgate active_outside_nwell = active - nwell nactive = active_outside_nwell & nplus ngate = nactive & poly nsd = nactive - ngate # Device extraction # PMOS transistor device extraction extract_devices(mos4("PMOS"), { "SD" => psd, "G" => pgate, "W" => nwell, "tS" => psd, "tD" => psd, "tG" => poly, "tW" => nwell }) # NMOS transistor device extraction extract_devices(mos4("NMOS"), { "SD" => nsd, "G" => ngate, "W" => bulk, "tS" => nsd, "tD" => nsd, "tG" => poly, "tW" => bulk }) # Define connectivity for netlist extraction # Inter-layer connect(psd, contact) connect(nsd, contact) connect(poly, contact) connect(contact, metal1) connect(metal1, metal1_lbl) # attaches labels connect(metal1, via1) connect(via1, metal2) connect(metal2, metal2_lbl) # attaches labels # Global connect_global(bulk, "SUBSTRATE") connect_global(nwell, "NWELL") # Compare section schematic("inv.cir") align # flattens unpaired circuits netlist.simplify # removes floating nets, combines devices compare
\end{lstlisting}

想试这段脚本，可以加载 KLayout 源码中的反相器版图 \ui{testdata/lvs/inv.oas}，打开宏编辑器 IDE（\menu{Tools > Macro Development}），在 LVS 标签页新建脚本并粘贴上面的内容，然后运行。LVS 报告浏览器会打开，并把一切显示为绿色。

\klfig[0.8\textwidth]{lvs_browser.png}{LVS 报告浏览器：全部匹配时显示为绿色}

\subsection{LVS 脚本的结构}

LVS 脚本中第一条、也是最重要的一条语句应是 \ui{deep} 开关，它启用层次化模式。若不启用层次化模式，生成的 netlist 不含 subcircuit；这种平坦 netlist 比对效率低，也难调试。因此用 \ui{deep} 切到层次化模式：

\begin{lstlisting}
deep
\end{lstlisting}

我们还指示 LVS 生成报告，并在结束后于报告浏览器中打开：

\begin{lstlisting}
report_lvs
\end{lstlisting}

需要的话也可以把报告写入文件（见 \ui{report\_lvs}）：

\begin{lstlisting}
report_lvs("inv.lvsdb")
\end{lstlisting}

接下来声明输入图层：

\begin{lstlisting}
nwell = input(1, 0) active = input(2, 0) pplus = input(3, 0) nplus = input(4, 0) poly = input(5, 0) contact = input(6, 0) metal1 = input(7, 0) metal1_lbl = labels(7, 1) via1 = input(8, 0) metal2 = input(9, 0) metal2_lbl = labels(9, 1)
\end{lstlisting}

\ui{input} 与 \ui{labels} 都是从版图源（如同 DRC，通常就是当前版图）取图层的函数。\ui{input} 取各类 Shape，而 \ui{labels} 只取 Text；这里用 \ui{labels} 取第一层金属的 label（GDS 第 7 层、datatype 1），第二层金属同理。

此外再创建一个空图层，用来表示 substrate。该层并不构成封闭区域，而是一堆 Shape，它们稍后会全部连到同一条（全局）net 上：

\begin{lstlisting}
bulk = polygon_layer
\end{lstlisting}

给图层起的名字其实就是代表版图图层的变量。和 DRC 一样，可以用它们算出若干派生图层：

\begin{lstlisting}
active_in_nwell = active & nwell pactive = active_in_nwell & pplus pgate = pactive & poly psd = pactive - pgate active_outside_nwell = active - nwell nactive = active_outside_nwell & nplus ngate = nactive & poly nsd = nactive - ngate
\end{lstlisting}

这些式子都是布尔运算：\ui{\&} 是逻辑与，\ui{-} 是逻辑非。因此 \ui{active\_in\_nwell} 是 ACTIVE 位于 NWELL 之内的部分，\ui{active\_outside\_nwell} 则是其外的部分。这些式子的主要目的是把 PMOS 与 NMOS 的源漏区区分开。

我们还分别取出 PMOS（pgate）与 NMOS（ngate）的栅区。准备好这些材料，就可以进入器件提取：

\begin{lstlisting}
extract_devices(mos4("PMOS"), { "SD" => psd, "G" => pgate, "W" => nwell, "tS" => psd, "tD" => psd, "tG" => poly, "tW" => nwell })
\end{lstlisting}

\ui{extract\_devices} 的第一个参数是器件提取器，即负责实际提取某类器件的对象。这里的模板是 \ui{MOS4}，而我们想产生一个名为 \ui{PMOS} 的新器件类，\ui{mos4("PMOS")} 即创建这样的提取器。

第二个参数是“图层符号 → 图层”的映射。每种器件提取器都定义了一套特定的图层符号；对任何器件都需要两类图层：提取器用来识别器件的输入图层，以及器件的 terminal 连接层。

输入图层用大写字母命名，terminal 输出图层则用小写 \ui{t} 加 terminal 名表示。上面那份声明是完整的，但由于 \ui{tW} 默认取 \ui{W}、\ui{tS} 与 \ui{tD} 默认取 \ui{SD}，可以写得更简短：

\begin{lstlisting}
extract_devices(mos4("PMOS"), { "SD" => psd, "G" => pgate, "W" => nwell, "tG" => poly })
\end{lstlisting}

\ui{NMOS} 类也需要一个提取器，构造方式与 PMOS 完全相同：

\begin{lstlisting}
extract_devices(mos4("NMOS"), { "SD" => nsd, "G" => ngate, "W" => bulk, "tS" => nsd, "tD" => nsd, "tG" => poly, "tW" => bulk })
\end{lstlisting}

有了器件，工作就完成了一半。接下来要交代连接关系（见 \ui{connect}）：

\begin{lstlisting}
connect(psd, contact) connect(nsd, contact) connect(poly, contact) connect(contact, metal1) connect(metal1, metal1_lbl) # attaches labels connect(metal1, via1) connect(via1, metal2) connect(metal2, metal2_lbl) # attaches labels
\end{lstlisting}

这些语句把 PMOS 源漏区（psd）与 CONTACT 相连，把 NMOS 源漏区（nsd）也与 CONTACT 相连，POLY 同样连到 CONTACT。注意我们在器件提取时已把 psd、nsd、poly 指定为 terminal 输出 \ui{tS}、\ui{tD}、\ui{tG}；把它们纳入连接关系之后，这些区域就成为对应 terminal 的一部分。

金属叠层很简单（CONTACT 接 METAL1，METAL1 经 VIA1 接 METAL2）。把 label 图层纳入连接关系，label 就自动挂到 net 上：net 提取器会从这些相连的 Text 对象中取出字符串，作为相应 net 的名字。

此外还需要做两处特殊连接（见 \ui{connect\_global}）：

\begin{lstlisting}
connect_global(bulk, "SUBSTRATE") connect_global(nwell, "NWELL")
\end{lstlisting}

全局连接的意思是：某图层上的所有 Shape 都属于同一条 net——即便彼此不接触——且这条 net 在 circuit 与 subcircuit 之间共享。对 bulk 层而言这当然是正确的，但对 NWELL 层未必：孤立的 NWELL 区块并不应连在一起。

至此，生成 netlist 所需的全部输入都已给定，接下来指定参考 netlist：

\begin{lstlisting}
schematic("inv.cir")
\end{lstlisting}

还有两步可选但值得推荐：层次对齐与提取网表的简化：

\begin{lstlisting}
align # flattens unpaired circuits netlist.simplify # removes floating nets, combines devices
\end{lstlisting}

\ui{align} 会把另一个 netlist 中不存在的 circuit 去掉，把其内容并入父 Cell。这样就能甩掉版图中常见、但 schematic 中并没有对应 Cell 的辅助单元（例如器件 PCell）。\ui{netlist.simplify} 则通过移除悬空 net、合并器件（例如指状晶体管）来精简 netlist，并会为接口建立 pin。

顺序应是先 \ui{align}，再 \ui{netlist.simplify}；两者都必须在 \ui{compare} 之前执行才有效。\ui{align} 见第 \ref{sec:lvs_overview} 节，\ui{netlist.simplify} 见第 \ref{sec:lvs_tweaks} 节。

最后，一切就绪后即可触发比对：

\begin{lstlisting}
compare
\end{lstlisting}

若在脚本开头插入一条写出 netlist 的语句（见 \ui{target\_netlist}），就能得到提取网表的 SPICE 版本：

\begin{lstlisting}
# SPICE output statement (insert at beginning of script): target_netlist("inv_extracted.cir", write_spice, "Extracted by KLayout")
\end{lstlisting}

由于 LVS 匹配成功，提取网表与参考网表几乎一致，只是多了源漏区面积、周长之类几何参数：

\begin{lstlisting}
* Extracted by KLayout * cell INVERTER .SUBCKT INVERTER * net 1 IN * net 2 VSS * net 3 VDD * net 4 OUT * net 5 NWELL * net 6 SUBSTRATE * device instance $1 r0 *1 1.025,4.95 PMOS M$1 3 1 4 5 PMOS L=0.25U W=1.5U AS=0.675P AD=0.675P PS=3.9U PD=3.9U * device instance $2 r0 *1 1.025,0.65 NMOS M$2 2 1 4 6 NMOS L=0.25U W=0.9U AS=0.405P AD=0.405P PS=2.7U PD=2.7U .ENDS INVERTER
\end{lstlisting}

\subsection{带下压二极管的反相器}

上面那个反相器 Cell 单独使用并没有太大用处，因为它缺少把 n-well 与 substrate 固定在确定电位的手段。这要靠下压二极管（tie-down diode）来实现。

下压二极管就是有源区上的接触。这些有源区在 substrate 上注 p+、在 n-well 内注 n+（与晶体管相反的掺杂类型）。有了这样的掺杂，金属接触就不会在硅体上形成肖特基势垒，而是表现为欧姆接触。所以严格说，这种“二极管”并不是整流意义上的真正二极管。

修改后的版图如下：

\klfig[0.8\textwidth]{inv_with_diodes.png}{带下压二极管的版图}

对应的 schematic 如下：

\klfig[0.8\textwidth]{inv_schematic2.png}{带下压二极管的 schematic}

有了这个电路，n-well 始终处于 VDD 电位，substrate 则被拉到 VSS：

\begin{lstlisting}
* Simple CMOS inverer circuit .SUBCKT INVERTER_WITH_DIODES VSS IN OUT VDD Mp VDD IN OUT VDD PMOS W=1.5U L=0.25U Mn OUT IN VSS VSS NMOS W=0.9U L=0.25U .ENDS
\end{lstlisting}

把下压二极管的提取包括进来之后，LVS 脚本略长一些：

\begin{lstlisting}
# LVS script (demo technology, KLayout manual) # Preamble: deep # Reports generated: report_lvs # LVS report window # Drawing layers: nwell = input(1, 0) active = input(2, 0) pplus = input(3, 0) nplus = input(4, 0) poly = input(5, 0) contact = input(6, 0) metal1 = input(7, 0) metal1_lbl = labels(7, 1) via1 = input(8, 0) metal2 = input(9, 0) metal2_lbl = labels(9, 1) # Bulk layer for terminal provisioning bulk = polygon_layer # Computed layers active_in_nwell = active & nwell pactive = active_in_nwell & pplus pgate = pactive & poly psd = pactive - pgate ntie = active_in_nwell & nplus active_outside_nwell = active - nwell nactive = active_outside_nwell & nplus ngate = nactive & poly nsd = nactive - ngate ptie = active_outside_nwell & pplus # Device extraction # PMOS transistor device extraction extract_devices(mos4("PMOS"), { "SD" => psd, "G" => pgate, "W" => nwell, "tS" => psd, "tD" => psd, "tG" => poly, "tW" => nwell }) # NMOS transistor device extraction extract_devices(mos4("NMOS"), { "SD" => nsd, "G" => ngate, "W" => bulk, "tS" => nsd, "tD" => nsd, "tG" => poly, "tW" => bulk }) # Define connectivity for netlist extraction # Inter-layer connect(psd, contact) connect(nsd, contact) connect(poly, contact) connect(ntie, contact) connect(nwell, ntie) connect(ptie, contact) connect(contact, metal1) connect(metal1, metal1_lbl) # attaches labels connect(metal1, via1) connect(via1, metal2) connect(metal2, metal2_lbl) # attaches labels # Global connect_global(bulk, "SUBSTRATE") connect_global(ptie, "SUBSTRATE") # Compare section schematic("inv2.cir") align netlist.simplify compare
\end{lstlisting}

主要差别在于计算 n 型下压区（n-well 内）与 p 型下压区的区域，这很简单：

\begin{lstlisting}
ntie = active_in_nwell & nplus ptie = active_outside_nwell & pplus
\end{lstlisting}

器件提取不变，但要把下压区纳入连接关系：

\begin{lstlisting}
connect(ntie, contact) connect(nwell, ntie) connect(ptie, contact)
\end{lstlisting}

把 ntie 连到 contact 与 nwell 之后，就连通到了 n-well，n-well 因而表现得像一个导电层（尽管电阻会很大）。别忘了 PMOS 的器件提取器把 bulk terminal 也放在 nwell 上，因此晶体管会自动连到 n-well 所在的 net。

ptie 不能简单地这样连，因为“substrate”并没有对应的 Polygon；但可以把 ptie 纳入全局连接：

\begin{lstlisting}
connect_global(bulk, "SUBSTRATE") connect_global(ptie, "SUBSTRATE")
\end{lstlisting}

此时 nwell 不再列入全局连接，于是不会再错误地把所有 nwell 区域都当作连在一起。

提取出的网表显示 NMOS 与 PMOS 的 bulk terminal 都接在源极上（漏源在此等价）：

\begin{lstlisting}
* Extracted by KLayout * cell INVERTER_WITH_DIODES .SUBCKT INVERTER_WITH_DIODES * net 1 IN * net 2 VDD * net 3 OUT * net 4 VSS * device instance $1 r0 *1 1.025,4.95 PMOS M$1 2 1 3 2 PMOS L=0.25U W=1.5U AS=0.675P AD=0.675P PS=3.9U PD=3.9U * device instance $2 r0 *1 1.025,0.65 NMOS M$2 4 1 3 4 NMOS L=0.25U W=0.9U AS=0.405P AD=0.405P PS=2.7U PD=2.7U .ENDS INVERTER_WITH_DIODES
\end{lstlisting}

\section{LVS 器件}\label{sec:lvs_devices}

\subsection{器件提取器与器件类}

KLayout 用两个概念来处理器件的多样性：

device class（器件类）即器件的分类。既有电阻、MOS 晶体管这类较宽泛的类别，也可以建立更具体的类别来表示具体形态，例如 NMOS 与 PMOS。器件类还决定器件如何合并。

器件类的文档见第 \ref{sec:lvs_device_classes} 节。

device extractor（器件提取器）才是真正分析版图并产出器件的对象。与器件类一样，提取器也有通用类型，每种提取器产出属于某个特定类的器件。

器件提取器的文档见第 \ref{sec:lvs_device_extractors} 节。

\section{LVS 器件类}\label{sec:lvs_device_classes}

KLayout 实现了多种标准器件类，它们是构造各类具体器件的基础。例如 \ui{MOS4} 类就是 NMOS 与 PMOS 晶体管具体器件类的基础。

\subsection{Resistor（电阻）}

\klfig[0.8\textwidth]{res_schematic.png}{普通电阻}

普通电阻有两个 terminal：A 与 B，具有以下参数：

\begin{kldesc}
\item R：电阻值，单位欧姆
\item L：长度，单位微米
\item W：宽度，单位微米
\item A：电阻区面积，单位平方微米
\item P：电阻区周长，单位微米
\end{kldesc}

电阻可以按并联或串联方式合并。

在 SPICE 中，普通电阻由 \ui{R} 元件表示；对应的 API 类见第 \ref{sec:lvs_device_classes} 节。

\subsection{带 bulk terminal 的电阻}

\klfig[0.8\textwidth]{res_with_bulk_schematic.png}{带 bulk terminal 的电阻}

带 bulk terminal 的电阻是普通电阻的扩展：参数相同，但多一个 terminal \ui{W}，连到电阻所在的区域（例如 well 或 substrate）。

若这些电阻的 bulk terminal 连在同一条 net 上，它们可以按并联或串联方式合并。

带 bulk 电阻的 API 类见第 \ref{sec:lvs_device_classes} 节。

\subsection{Capacitor（电容）}

\klfig[0.8\textwidth]{cap_schematic.png}{普通电容}

普通电容有两个 terminal：A 与 B，具有以下参数：

\begin{kldesc}
\item C：电容值，单位法拉
\item A：电容区面积，单位平方微米
\item P：电容区周长，单位微米
\end{kldesc}

在 SPICE 中，普通电容由 \ui{C} 元件表示；对应的 API 类见第 \ref{sec:lvs_device_classes} 节。

\subsection{带 bulk terminal 的电容}

\klfig[0.8\textwidth]{cap_with_bulk_schematic.png}{带 bulk terminal 的电容}

带 bulk terminal 的电容是普通电容的扩展：参数相同，但多一个 terminal \ui{W}，连到电容所在的区域（例如 well 或 substrate）。

若这些电容的 bulk terminal 连在同一条 net 上，它们可以按并联或串联方式合并。

带 bulk 电容的 API 类见第 \ref{sec:lvs_device_classes} 节。

\subsection{Diode（二极管）}

\klfig[0.8\textwidth]{diode_schematic.png}{二极管}

二极管有两个 terminal：A 与 C，分别为阳极与阴极，具有以下参数：

\begin{kldesc}
\item A：二极管面积，单位平方微米
\item P：二极管周长，单位微米
\end{kldesc}

二极管按并联方式合并（A 对 A、C 对 C），合并后其面积与周长相加。

在 SPICE 中，二极管由 \ui{D} 元件表示，并以器件类名作为 model 名；对应的 API 类见第 \ref{sec:lvs_device_classes} 节。

\subsection{MOS 晶体管}

\klfig[0.8\textwidth]{mos3_schematic.png}{三端 MOS 晶体管}

三端 MOS 晶体管的 terminal 为 S、G、D，即源、栅、漏，其中 S 与 D 可互换。其参数为：

\begin{kldesc}
\item W：栅宽，单位微米
\item L：栅（沟道）长，单位微米
\item AS：源区面积，单位平方微米
\item PS：源区周长，单位微米
\item AD：漏区面积，单位平方微米
\item PD：漏区周长，单位微米
\end{kldesc}

当源漏相连、栅极相连且栅长完全相同时，MOS3 晶体管按并联方式合并，合并后宽度、面积与周长相加。

\subsection{带 bulk 的 MOS 晶体管}

三端 MOS 晶体管的 API 类见第 \ref{sec:lvs_device_classes} 节。

\klfig[0.8\textwidth]{mos4_schematic.png}{四端 MOS 晶体管}

四端晶体管是三端晶体管的扩展，多一个 bulk terminal \ui{B}。这大概是最重要的晶体管器件，因为四端版本与 SPICE 的 \ui{M} 元件兼容。

若这些 MOS 晶体管的 bulk terminal 连在同一条 net 上，它们可以像三端版本那样按并联方式合并。

在 SPICE 中，MOS4 器件由 \ui{M} 元件表示，并以器件类名作为 model 名；对应的 API 类见第 \ref{sec:lvs_device_classes} 节。

\subsection{双极晶体管}

\klfig[0.8\textwidth]{bjt3_schematic.png}{三端双极晶体管}

三端双极晶体管可以是 NPN 型或 PNP 型。在 KLayout 中，这种器件类型既能表示横向型也能表示纵向型。其参数为：

\begin{kldesc}
\item AE：发射区面积，单位平方微米
\item PE：发射区周长，单位微米
\item NE：发射极个数（初始为 1）
\item AB：基区面积，单位平方微米
\item PB：基区周长，单位微米
\item AC：集电区面积，单位平方微米
\item PC：集电区周长，单位微米
\end{kldesc}

提取时，多发射极版本会被提取成多个器件——每个发射区一个——此时 NE = 1。双极晶体管按并联方式合并，合并后发射极参数 AE、PE 与 NE 相加。

在 SPICE 中，BJT3 器件由 \ui{Q} 元件表示，并以器件类名作为 model 名；对应的 API 类见第 \ref{sec:lvs_device_classes} 节。

\subsection{带衬底的双极晶体管}

\klfig[0.8\textwidth]{bjt4_schematic.png}{带衬底的四端双极晶体管}

四端晶体管是三端晶体管的扩展，多一个 bulk terminal \ui{S}。

若这些双极晶体管的 bulk terminal 连在同一条 net 上，它们可以像三端版本那样按并联方式合并。

在 SPICE 中，BJT4 器件由带四个节点的 \ui{Q} 元件表示，并以器件类名作为 model 名；对应的 API 类见第 \ref{sec:lvs_device_classes} 节。

\section{LVS 器件提取器}\label{sec:lvs_device_extractors}

器件提取器是器件提取过程中真正的“干活者”。KLayout 自带多种现成的器件提取器，也可以在 LVS 脚本框架内（用 Ruby 编写）实现自定义提取器。

\subsection{电阻提取器（resistor 与 resistor\textbackslash{}\_with\textbackslash{}\_bulk）}

电阻提取器假设版图由一段电阻“连线”加两个端帽（contact）组成：连线用图层符号 \ui{R} 指定，端帽用 \ui{C} 指定。

\klfig[0.8\textwidth]{res_ex_layout.png}{电阻提取器的输入：R 连线与 C 端帽}

提取器由方块数（number of squares）与方块电阻（sheet resistance）计算阻值，方块电阻在创建提取器时给出。

\begin{lstlisting}
sheet_rho = 0.5 model_name = "RES" extract_devices(resistor(model_name, sheet_rho), { "R" => res_layer, "C" => contact_layer })
\end{lstlisting}

普通电阻有两个 terminal，分别输出到 \ui{tA} 与 \ui{tB} 两个 terminal 图层上。若未指定 \ui{tA} 或 \ui{tB}，则 \ui{A} 或 \ui{B} terminal 分别写到 \ui{C} 图层上。

\klfigpair{res_ex_ta.png}{tA terminal 输出}{res_ex_tb.png}{tB terminal 输出}

对带 bulk 的电阻，连线区域作为 \ui{W} terminal 输出到 \ui{tW} terminal 图层。

\klfig[0.8\textwidth]{res_ex_tw.png}{tW terminal 输出}

注意：阻值计算用的是一种简单近似——沿 \ui{R} Polygon 的周长走一圈，把周长分成与 \ui{C} 层接触的部分和不接触的部分，由此得出方块数。因此这是一种粗略但通常可用的估算。

\subsection{电容提取器（capacitor 与 capacitor\textbackslash{}\_with\textbackslash{}\_bulk）}

电容被假定由两块“极板”构成（即纵向电容），极板分别在 \ui{P1} 与 \ui{P2} 层上；电容取自这两层的重叠区域。

\klfig[0.8\textwidth]{cap_ex_layout.png}{纵向电容：两块极板 P1 与 P2 的重叠}

提取器由重叠面积与单位面积电容（F/平方微米）计算电容值。

\begin{lstlisting}
area_cap = 1.5e-15 model_name = "CAP" extract_devices(capacitor(model_name, area_cap), { "P1" => metal1, "P2" => metal2 })
\end{lstlisting}

普通电容有两个 terminal，分别输出到 \ui{tA} 与 \ui{tB} 图层。若未指定，\ui{A} 与 \ui{B} terminal 分别写到 \ui{P1} 与 \ui{P2} 图层上。

\klfigpair{cap_ex_ta.png}{tA terminal}{cap_ex_tb.png}{tB terminal}

对带 bulk 的电容，电容区域作为 \ui{W} terminal 输出到 \ui{tW} terminal 图层。

\klfig[0.8\textwidth]{cap_ex_tw.png}{tW terminal}

\subsection{二极管提取器（diode）}

二极管被假定由两个纵向注入区（well、扩散区）构成：一个是 p 型（\ui{P} 层），另一个是 n 型（\ui{N} 层）。这两层同时也构成二极管的阳极（p）与阴极（n）。

\klfig[0.8\textwidth]{diode_ex_layout.png}{二极管提取输入：P 与 N 两个注入区}

\begin{lstlisting}
model_name = "DIODE" extract_devices(diode(model_name), { "P" => pplus, "N" => nwell })
\end{lstlisting}

二极管有两个 terminal，分别输出到 \ui{tA} 与 \ui{tC} 图层。若未指定 \ui{tA}，\ui{A} terminal 写到 \ui{P} 层；若未指定 \ui{tC}，\ui{C} terminal 写到 \ui{N} 层。

\klfigpair{diode_ex_ta.png}{tA terminal}{diode_ex_tc.png}{tC terminal}

\subsection{MOS 晶体管提取器（mos3 与 mos4）}

MOS 晶体管靠栅极（\ui{G} 输入）与源漏区（\ui{SD} 输入）识别：源漏区必须与栅极图形分离，栅极与源漏区相接触的边决定器件宽度，而垂直于该边的方向决定沟道长度。

\klfig[0.8\textwidth]{mos_ex_layout.png}{MOS 提取输入：G 与 SD}

\begin{lstlisting}
model_name = "PMOS" extract_devices(mos4(model_name), { "SD" => (active - poly) & pplus, "G" => (active & poly), "W" => nwell, "tS" => active, "tD" => active, "tG" => poly, "tB" => nwell })
\end{lstlisting}

MOS3 器件产生三个 terminal，分别输出到 \ui{tS}、\ui{tG}、\ui{tD} 图层（依次为源、栅、漏）。

\klfigtriple{mos_ex_ts.png}{tS terminal}{mos_ex_tg.png}{tG terminal}{mos_ex_td.png}{tD terminal}

MOS4 器件多一个 terminal（bulk），写到 \ui{tB}。

\klfig[0.8\textwidth]{mos_ex_tb.png}{MOS4 的 bulk terminal}

\subsection{扩散 MOS 晶体管提取器（dmos3 与 dmos4）}

DMOS 器件与 MOS 器件基本相同，区别在于其源、漏是区分开的。扩散型 MOS 晶体管常常如此，故得此名。

DMOS 与 MOS 器件共用同一个器件类。DMOS 器件被配置为源漏不可互换，因此若它们被互换，netlist 比对会报告为错误。

DMOS 晶体管靠栅极（\ui{G} 输入）、源区（\ui{S} 输入）与漏区（\ui{D} 输入）识别；源漏区须与栅极分离，其与栅极相接触的边决定器件宽度，垂直方向的尺寸决定沟道长度。

\subsection{双极晶体管提取器（bjt3 与 bjt4）}

双极晶体管基本上分两类：纵向型与横向型。

横向晶体管由注入或扩散的 well 在晶圆上形成交替的 n/p 结构，其基本识别区是基区：集电区与发射区位于基区之内或与其重叠，采用相反的掺杂类型。

\klfig[0.8\textwidth]{bjt_lateral.png}{横向双极晶体管的结构}

纵向晶体管由叠放的 n/p well 构成，有时在标准 CMOS 工艺中会作为寄生成分出现。例如以 substrate 作集电区、nwell 作基区、pplus 注入作发射区，就能构成一个 PNP 晶体管。

\klfig[0.8\textwidth]{bjt_vertical.png}{纵向双极晶体管的结构}

两种情况下，一个基区之内都可能有多个发射区；此时每个发射区各提取一个晶体管。

\subsubsection{纵向双极晶体管}

纵向双极晶体管的输入取自 \ui{B}（基区）、\ui{C}（集电区）与 \ui{E}（发射区），其中 \ui{C} 是可选的。

\klfig[0.8\textwidth]{bjt_ex_layout.png}{纵向双极：输入 B、C、E}

尤其对双极器件而言，给出有用的 terminal 输出层很重要。通常 well 与扩散区会通过 contact 相连（此处暂不考虑肖特基二极管），因此把 terminal 送到与 contact 相接触的形状上是明智的做法。

\begin{lstlisting}
model_name = "PNP" extract_devices(bjt3(model_name), { "C" => collector, "B" => base, "E" => emitter, "tC" => contact, "tB" => contact, "tE" => contact })
\end{lstlisting}

BJT3 器件产生三个 terminal，分别输出到 \ui{tC}、\ui{tB}、\ui{tE} 图层（依次为集电、基区、发射区）。

\klfigtriple{bjt_ex_tc.png}{tC terminal}{bjt_ex_tb.png}{tB terminal}{bjt_ex_te.png}{tE terminal}

若集电区为空（例如 p 型 substrate），则把基区形状复制到 \ui{tC} 输出层作为集电区 terminal。

BJT4 器件多一个 terminal（substrate），写到 \ui{tS}。\ui{tS} 是发射区形状的副本，但连到 substrate terminal 上。

\klfig[0.8\textwidth]{bjt_ex_ts.png}{BJT4 的 substrate terminal}

\subsubsection{横向双极晶体管}

横向双极晶体管的输入同样取自 \ui{B}、\ui{C} 与 \ui{E}。对横向晶体管，\ui{C} 不可省略，且不得完全覆盖基区。除此之外，BJT3 与 BJT4 提取器的用法完全相同。

典型的横向晶体管由基区内的一圈集电区（collector ring）与一个发射岛构成：

\klfig[0.8\textwidth]{bjtlat_ex_layout.png}{横向双极：集电环与发射岛}

横向情形下双极提取器产生的 terminal 与纵向情形相同，只是几何不同：

\klfigtriple{bjtlat_ex_tc.png}{横向情形的 terminal 几何}{bjtlat_ex_tb.png}{横向 BJT4 的 substrate terminal}{bjtlat_ex_te.png}{横向情形 terminal 输出}

同样地，对 BJT4，\ui{tS} 是发射区形状的副本但连到 substrate terminal。

\klfig[0.8\textwidth]{bjtlat_ex_ts.png}{横向 BJT4 的 tS}

\subsection{器件提取器与器件类}

\ui{extract\_devices} 会返回所生成器件的 DeviceClass 对象，可用它做一些基本调整，其中最重要的是启用或禁用某些参数。

多数器件提取器实际提取的参数比默认给出的更多。例如电阻提取器不只是求阻值，还会求出电阻条的长度 L、宽度 W 以及面积 A 与周长 P；默认只有其中一部分参与比对。

参数可以用 \ui{enable\_parameter} 完全启用，或用 \ui{disable\_parameter} 禁用；\ui{tolerance} 可以启用某参数参与比对并指定比对容差；\ui{ignore\_parameter} 则让某参数在比对步骤中被忽略。

定制内置器件提取器的另一种方式是提供自定义的器件类。下面的代码新建一个电阻类，修改预配置的器件参数定义以启用 \ui{W} 与 \ui{L}。

\begin{lstlisting}
class MyResistor < RBA::DeviceClassResistor def initialize super enable_parameter("W", true) enable_parameter("L", true) end end ... extract_devices(resistor("RES", 1, MyResistor), ...)
\end{lstlisting}

使用自定义器件类还带来一个可能：补充额外参数，或者借助电阻提取器的提取机制实现一种全新的器件。唯一的要求是参数与 terminal 定义必须兼容。

\section{LVS 输入输出}\label{sec:lvs_io}

LVS（就 netlist 提取而言 DRC 也一样）提供读写 netlist、schematic、标注版图与 LVS 结果的接口。I/O 分三大类：

\begin{itemize}
\item Netlist：纯粹的电路信息，配合 subcircuit 就构成层次化 netlist。目前可导入导出的格式是某种 SPICE netlist 方言，读写过程可以定制以获得一定灵活性。
\item Layout-to-netlist database（L2N DB）：也称 extracted netlist 或 annotated layout，即取自原版图的 netlist 连同相应的 Shape。凡涉及非器件的 Shape，该数据库都能在几何上重建 net；器件则被抽象成类似“黑盒”的表示。
\item LVS result database（LVS DB）：即 L2N 数据库加上参考 netlist 与 cross reference——后者是配对成功的 circuit、net、device、pin、subcircuit 及其状态信息的列表，既是查找表也是调试工具。
\end{itemize}

\subsection{写出 netlist}

可以把 netlist 写成文件，例如提供给（功能）仿真器。在 LVS 脚本中，全局语句 \ui{target\_netlist} 触发 netlist 写出（细节见 \ui{target\_netlist}）。

\begin{lstlisting}
target_netlist("output.cir", write_spice, "Created by KLayout")
\end{lstlisting}

这条语句原则上可以出现在脚本任何位置，netlist 会在脚本成功执行后写出。第一个参数是文件路径（默认相对于原版图文件），第二个参数是 writer：\ui{write\_spice} 创建一个 SPICE netlist writer。

\ui{write\_spice} 配置函数有两个选项：

\begin{lstlisting}
write_spice(use_net_names, with_comments)
\end{lstlisting}

两个选项都是布尔值。若存在且为真，第一个选项让 writer 使用真实的 net 名而不是数字编号；\ui{with\_comments} 若存在且为真，会在 netlist 中嵌入调试注释，显示实例位置、pin 名等。

还可以通过给 SPICE writer 提供 delegate 来做更多定制。做法是继承 \kte{NetlistSpiceWriterDelegate} 并重实现其中一个或多个预留方法。

下面这个例子不输出器件元件，而是输出 subcircuit 调用：

\begin{lstlisting}
# Write extracted netlist to extracted.cir using a special # writer delegate # This delegate makes the writer emit subcicuit calls instead of # standard elements for the devices class SubcircuitModels < RBA::NetlistSpiceWriterDelegate def write_header emit_line(".INCLUDE 'models.cir'") end def write_device(device) str = "X" + device.expanded_name device_class = device.device_class device_class.terminal_definitions.each do |td| str += " " + net_to_string(device.net_for_terminal(td.id)) end str += " " + device_class.name str += " PARAMS:" device_class.parameter_definitions.each do |pd| str += " " + pd.name + ("=%.12g" % device.parameter(pd.id)) end emit_line(str) end end # Prepare a writer using the new delegate custom_spice_writer = RBA::NetlistSpiceWriter::new(SubcircuitModels::new) custom_spice_writer.use_net_names= true custom_spice_writer.with_comments = false # The declaration of netlist production using the new custom writer target_netlist("extracted.cir", custom_spice_writer, "Extracted by KLayout")
\end{lstlisting}

对 LVS 入门一节中的简单反相器，这段脚本会产生如下 netlist。它不再打印默认的 \ui{M} 元件，而是产生 subcircuit 调用，从而可以在 subcircuit 里放更精细的模型；器件类名就用来指向这些模型 subcircuit。

\begin{lstlisting}
* Extracted by KLayout .INCLUDE 'models.cir' .SUBCKT INVERTER X$1 VDD IN OUT NWELL PMOS PARAMS: L=0.25 W=1.5 AS=0.675 AD=0.675 PS=3.9 PD=3.9 X$2 VSS IN OUT SUBSTRATE NMOS PARAMS: L=0.25 W=0.9 AS=0.405 AD=0.405 PS=2.7 + PD=2.7 .ENDS INVERTER
\end{lstlisting}

netlist 也可以直接从 netlist 对象写出：脚本中先用 \ui{netlist} 函数取得对象（该函数若尚未提取会先触发一次提取），再调用它的写出方法。与 \ui{target\_netlist} 不同，这种方式在调用处立即写出，而不是等脚本结束。

\subsection{读取 netlist}

读取 netlist 的主要用途是在 LVS 中作比对：参考 netlist 用 \ui{schematic} 函数读入（见 \ui{schematic}）。

\begin{lstlisting}
schematic("inverter.cir")
\end{lstlisting}

目前的 SPICE 支持有一些限制：默认只实现了部分元件，即 \ui{M}（得到 \ui{MOS4} 器件类）、\ui{Q}（得到 BJT3 或 BJT4）、\ui{R}（电阻）、\ui{C}（电容）与 \ui{D}（二极管）。

与 SPICE writer 一样，reader 也可以提供 delegate 来定制：继承相应类并重实现其预留方法即可。

下面这个例子定制 reader，从 subcircuit 模型而不是 \ui{M} 元件取 MOS 器件；它基本上正是前面 writer 定制的反向操作（仅针对 MOS 器件）。

\begin{lstlisting}
# Provides a SPICE netlist reader delegate which turns # some subcircuit models (for subcircuits NMOS and PMOS) # into devices class SubcircuitModelsReader < RBA::NetlistSpiceReaderDelegate # implements the delegate interface: # says we want to catch these subcircuits as devices def wants_subcircuit(name) name == "NMOS" || name == "PMOS" end # implements the delegate interface: # take and translate the element def element(circuit, el, name, model, value, nets, params) if el != "X" # all other elements are left to the standard implementation return super end if nets.size != 4 error("Subcircuit #{model} needs four nodes") end # provide a device class cls = circuit.netlist.device_class_by_name(model) if ! cls cls = RBA::DeviceClassMOS4Transistor::new cls.name = model circuit.netlist.add(cls) end # create a device device = circuit.create_device(cls, name) # and configure the device [ "S", "G", "D", "B" ].each_with_index do |t,index| device.connect_terminal(t, nets[index]) end # parameters in the model are given in micrometer units, so # we need to translate the parameter values from SI to um values: device.set_parameter("W", (params["W"] || 0.0) * 1e6) device.set_parameter("L", (params["L"] || 0.0) * 1e6) return true end end # Instantiate a reader using the new delegate reader = RBA::NetlistSpiceReader::new(SubcircuitModelsReader::new) # Import the schematic with this reader schematic("inv_xmodels.cir", reader)
\end{lstlisting}

\subsection{Layout-to-Netlist 数据库 / 报告}

layout-to-netlist 数据库（L2N DB）用全局函数 \ui{report\_netlist} 写出。该函数可放在脚本任意位置，写出发生在脚本成功执行之后：

\begin{lstlisting}
report_netlist("extracted.l2n")
\end{lstlisting}

若不给文件名，则只打开 netlist 浏览器而不写文件。layout-to-netlist 数据库是 KLayout 专用格式，含 netlist 信息以及版图中的 Shape 与 Instance 信息；L2N 数据库可以读入 netlist 浏览器中查看。

\subsection{Layout-vs-Schematic 数据库 / 报告}

Layout-vs-schematic 数据库（LVS DB）用全局函数 \ui{report\_lvs} 写出。该函数同样可放在脚本任意位置，写出发生在脚本成功执行之后：

\begin{lstlisting}
report_lvs("extracted.lvsdb")
\end{lstlisting}

若不给文件名，则只打开 netlist 浏览器而不写文件。LVS 数据库是 KLayout 专用格式，含提取网表信息、参考网表以及 cross reference 表；LVS 数据库可以读入 netlist 浏览器中查看。

\section{LVS 连接关系}\label{sec:lvs_connect}

\subsection{层内连接与层间连接}

LVS 脚本中的连接关系设置决定了连接如何建立。连接通常通过铝、铜这类导电材料形成：表示同一材料的 Polygon 相互连通；连接也可以跨越多个 Polygon，只要它们彼此接触即构成同一导电区域。

连接常常跨层进行，例如 via 就是绝缘层上的一个孔，把两层金属连起来。这类连接用 \ui{connect} 语句建模（见 \ui{connect}）：

\begin{lstlisting}
connect(layer1, layer2)
\end{lstlisting}

connect 语句表示：当 layer1 与 layer2 的 Polygon 相互重叠时，二者电气相连。layer1 与 layer2 可以是原始图层或派生图层。\ui{connect} 语句应写在需要 netlist 之前——也就是 \ui{compare} 或任何读取 netlist 的语句之前。

下面是一个三层金属叠层的纵向剖面示例，以及相应的 \ui{connect} 语句：

\klfig[0.8\textwidth]{metal_connections.png}{}

label 也可以纳入连接关系。label 通常放在金属层上：若 label 与金属 Shape 画在同一层，用 \ui{input} 读该层时它们会自动被包含；若只想从某层读 label，用 \ui{labels} 函数。

要把 label 挂到金属层上，只需把 label 层与金属层连接起来：

\begin{lstlisting}
metal1_labels = labels(10, 0) metal1 = input(11, 0) via1 = input(12, 0) metal2_labels = labels(13, 0) metal2 = input(14, 0) connect(metal1, metal1_labels) connect(metal1, via1) connect(via1, metal2) connect(metal2, metal2_labels)
\end{lstlisting}

label 连到金属层之后，其文字将用于给形成的 net 命名。理想情况下，一条 net 上只有一段文字，或有多段但内容相同的文字，此时 net 名没有歧义；若一条 net 上有多个内容不同的 label，KLayout 会做出选择并给出警告。

\subsection{全局连接}

KLayout 支持“隐式全局连接”：某图层上的所有 Polygon 一律视为相连，无论它们在几何上是否接触。典型例子是 substrate（即 bulk），它代表略有导电性的衬底材料，而版图中并没有表示整块衬底的 Polygon。

\begin{lstlisting}
connect_global(bulk, "VSS")
\end{lstlisting}

\ui{connect\_global} 的参数是全局连接的图层以及要创建的全局 net 名。下面的语句把 \ui{bulk} 上的所有 Shape 连成一条 net，名为 \ui{VSS}：每个 circuit 至少都有 \ui{VSS} 这条 net。

\subsection{隐式连接}

隐式连接适合用来提供那些打算在更高层才建立的连接。设想某个 circuit 有一条很大的电源网：做版图时电源网可能还没有完全连通，因为计划在上层或由布线工具把它接起来。

为了能验证这样的 Cell，可以用“隐式连接”：给各部分 net 打上相同名字的 label，并假定它们相连。例如要给 \ui{VDD} net 的各部分建立隐式连接，就在每部分上放一个 \ui{VDD} label，然后加上：

\begin{lstlisting}
connect_implicit("VDD")
\end{lstlisting}

\ui{connect\_implicit}（见 \ui{connect\_implicit}）可以出现多次，用来建立多条这样的连接。隐式连接在顶层被接受，但会给出警告，提示该连接应在更高层完成；在 subcircuit 中，隐式连接同样会被报出，除非后续得到物理连接的确认。

也可以把某个版图声明为顶层版图，这会把“缺少物理连接”的警告升级为错误：

\begin{lstlisting}
top_level(true)
\end{lstlisting}

在其他系统里，\ui{connect\_implicit} 这类功能也被称为“must connect”网。

可以用 glob 风格的通配符，把某一类 label 纳入 \ui{connect\_implicit} 语句：

\begin{lstlisting}
connect_implicit("VDD*")
\end{lstlisting}

这样一来，所有带 \ui{VDD1} label 的 net（或带 \ui{VDD\_5V} label 的 net）都会各自连在一起。但该语句只把 \ui{VDD1} 与 \ui{VDD1} 相连，不会把 \ui{VDD1} 与 \ui{VDD2} 相连；要让不同名字的 label 之间建立连接，需另外写语句。

\ui{connect\_implicit} 可以出现多次，每个语句都在扩大“哪些 label 会被连接”的范围。

上面的写法适用于所有 Cell。也可以把它限定到具体 Cell：第一个参数给 Cell 名的 glob 通配符，第二个参数给 net 名通配符。

下面这条语句把 \ui{MEMCELL} 子 Cell 中所有带 \ui{VDD} label 的 net 连起来：

\begin{lstlisting}
connect_implicit("MEMCELL", "VDD")
\end{lstlisting}

该规则作用于所有匹配第一个参数通配符的 Cell。同样，\ui{connect\_implicit} 规则可以给多次，此时所有匹配到的规则共同起作用。

\ui{connect\_implicit} 语句必须写在 netlist 提取之前，通常在 \ui{connect} 语句之前或紧随其后。

\subsection{显式连接}

显式连接适合用来强制建立那种 schematic 中存在、但在该 Cell 层面版图上并未物理建立的连接。例如下面这个反相器版图：

\klfig[0.8\textwidth]{inv.png}{}

版图中没有下压二极管，因此 n-well 区没有物理连接，bulk 衬底也没有物理连接。这样能省面积，但这些二极管迟早要通过别的方式补上，通常是在把标准单元组合起来的时候。

然而在反相器这一层并没有这种连接，于是反相器有独立的 bulk 与 n-well pin；而 schematic 有时是简化版，不提供这些 pin。这样就出现了版图与 schematic 之间的固有失配。

\klfig[0.8\textwidth]{inv_explicit.png}{}

为使版图与 schematic 对齐，可以显式地把 bulk 与 VSS pin 相连，n-well 与 VDD 同理。这套机制与前面讲的 \ui{connect\_implicit} 类似，区别在于它能连接名字不同的 net。

要在上面的例子里建立显式连接，先确保 n-well 与 bulk 有合适的名字：n-well 通过在 n-well 岛上放 label 命名（例如 \ui{NWELL}）；bulk 并不是真正带 Polygon 的版图图层，因此用 \ui{connect\_global} 给它一个全局名。

下面的代码把 \ui{INV} Cell 内的 bulk net 与 \ui{VSS} 连起来：

\begin{lstlisting}
connect_global(bulk, "BULK") ... connect_explicit("INV", [ "BULK", "VSS" ])
\end{lstlisting}

注意该规则会形成一条名为 \ui{BULK,VSS} 的新 net，把两条子 net 合并。

Cell 名可以是通配符，例如 \ui{INV*} 会把该规则应用到所有以 \ui{INV} 开头的 Cell。Cell 通配符不是必须的，省略时规则作用于所有 Cell。

与隐式连接一样，显式连接也会检查其在上一层是否落实——无论是物理连接，还是另一条显式/隐式连接。

显式连接还会对所涉及 net 名隐含地建立隐式连接，因此上例中 \ui{VSS} 的各个片段即便物理上未相连也被视为相连，同样会检查这些连接是否在更高层落实。

同样，\ui{connect\_explicit} 语句必须写在 netlist 提取之前，通常在 \ui{connect} 语句之前或紧随其后。

\subsection{软连接}

软连接（soft connection）用来找出“信号甚至电源经由高阻通路走线”这类布线问题。高阻通路可以是经多晶硅、注入区、well 或衬底形成的连接，这类区域的电阻常比金属高几个数量级。

举一例：

\klfig[0.8\textwidth]{soft_connections.png}{}

这是教科书式的平面 CMOS 工艺，两个 PMOS 器件位于同一个 n-well 中，例如可分属两组反相器。两个 PMOS 的源极都需要接到 VDD，n-well 区域同样需要接到 VDD。

这样的工艺叠层可以这样描述连接关系：

\begin{lstlisting}
# Input layers nwell = ... active = ... pplus = ... nplus = ... poly = ... contact = ... metal1 = ... via1 = ... metal2 = ... # computed layers (nactive, pactive) = active.and_not(nwell) # PMOS and NMOS source/drain regions psd = (nactive & pplus) - poly nsd = (pactive & nplus) - poly # n tie and p tie (nwell and substrate contact) ntie = nactive & nplus ptie = pactive & pplus # connections # nwell connections connect(ntie, nwell) connect(contact, ntie) # substrate connections connect_global(ptie, "BULK") connect(contact, ptie) # device connections connect(contact, psd) connect(contact, nsd) connect(contact, poly) # metal connections connect(metal1, contact) connect(via1, metal1) connect(metal2, via1)
\end{lstlisting}

但问题来了：如图所示，左侧 PMOS 的源极确实接到了 VDD，右侧那个却缺少通往 VDD 的金属连接。仅从纯连通性看，这个晶体管的源极仍然连到了 VDD——只是经由了高阻的 n-well。

解决办法是引入软连接：把 \ui{connect} 换成 \ui{soft\_connect}，把 \ui{connect\_global} 换成 \ui{soft\_connect\_global}（见 \ui{soft\_connect} 与 \ui{soft\_connect\_global}）。第一个图层是“上层”，第二个是“下层”。

软连接可以理解为有方向的连接：电流只能从上层流向下层，反之不行。因此只有当软连接的两个上层 terminal 都接到同一条物理 net 上时，才构成真正的连接。

要解决 n-well 的问题，把 n-tie 到 n-well 的那条连接语句改成 \ui{soft\_connect}：

\begin{lstlisting}
soft_connect(ntie, nwell)
\end{lstlisting}

此时上面的图就变成这样：

\klfig[0.92\textwidth]{soft_connections2.png}{}

有了这个定义，netlist 提取器就能发现该故障并给出警告或错误（顶层模式下为错误）。警告显示在日志标签页，指出布线不完整，并给出被软连接分开的各子 net 的细节。

为完整起见，还需要纳入其它软连接，例如经衬底（到 BULK net 的全局软连接）、经源漏注入区、经 tie 注入区以及经多晶硅的连接：

\begin{lstlisting}
# Input layers nwell = ... active = ... pplus = ... nplus = ... poly = ... contact = ... metal1 = ... via1 = ... metal2 = ... # computed layers (nactive, pactive) = active.and_not(nwell) # PMOS and NMOS source/drain regions psd = (nactive & pplus) - poly nsd = (pactive & nplus) - poly # n tie and p tie (nwell and substrate contact) ntie = nactive & nplus ptie = pactive & pplus # connections # nwell connections soft_connect(ntie, nwell) soft_connect(contact, ntie) # substrate connections soft_connect_global(ptie, "BULK") soft_connect(contact, ptie) # device connections soft_connect(contact, psd) soft_connect(contact, nsd) soft_connect(contact, poly) # metal connections connect(metal1, contact) connect(via1, metal1) connect(metal2, via1)
\end{lstlisting}

如代码所示，可以指定多条软连接；从检查的角度看，所有软连接都是同一种性质。

注意，两条方向相反的软连接会相互抵消，从而变成一条硬连接：

\begin{lstlisting}
soft_connect(a, b) soft_connect(b, a)
\end{lstlisting}

注意：因此务必留意软连接的方向——上层先写高导电（低阻）图层，下层后写低导电（高阻）图层。

\subsubsection{软连接与“must connect”网}

软连接与 must connect 网（即 \ui{connect\_explicit} 与 \ui{connect\_implicit}）目的相同：发现不完整的布线。二者通常配合使用，例如对 VDD net 做 \ui{connect\_implicit}，并在 schematic 未把 VDD 与 NWELL 相连时对 VDD 与 NWELL 做 \ui{connect\_explicit}。

软连接在执行 \ui{connect\_explicit} 与 \ui{connect\_implicit} 之前就被检查，因此软连接错误无法通过声明 must connect 来掩盖。反过来这也不成问题，因为两种检查都会给出警告或错误（顶层模式下）。

\section{LVS 比对}\label{sec:lvs_compare}

实际的比对步骤相当简单：只要已经设置好提取（\ui{extract\_devices}）、连接关系（\ui{connect}、\ui{connect\_global}、\ui{connect\_implicit}），并提供了参考 netlist，就可以直接调用 \ui{compare}。

\begin{lstlisting}
compare
\end{lstlisting}

该方法（\ui{compare}）会提取 netlist（若尚未提取）并与 schematic 比对，成功返回 true，否则返回 false——若你想据此做进一步处理会很有用。

比对步骤可以通过设置提示（hint）来配置。

\subsection{net 等同提示}

把两条 net 声明为等同有时很有用，至少在调试阶段如此：比对算法随后可以据此推断失配的真正原因。例如接线名与 schematic 中的名字不一致时，这样可以排除干扰。

\begin{lstlisting}
same_nets("LOGIC", "VDD", "VDD:P")
\end{lstlisting}

这个例子假定电源网在版图中标为 \ui{VDD}、在 schematic 中叫 \ui{VDD:P}。正式验证时不要把这条语句留在脚本里。

\ui{same\_nets} 也可以用来要求指定的 net 必须匹配。这在顶层很有用，例如检查每个 pad 上分配的 net 是否对应，从而保证 pad 接对。

\begin{lstlisting}
same_nets!("CHIP", "VDD", "VDD")
\end{lstlisting}

带感叹号的写法含义不同：若版图或 schematic 中缺少 \ui{VDD} 网，或两者不匹配，就报告 net 失配。上面的写法还可以进一步组合使用。

\begin{lstlisting}
same_nets!("CHIP", "VDD")
\end{lstlisting}

circuit 名与 net 名都支持通配符。下面的例子要求 \ui{INVERTER} circuit 中所有以 \ui{PAD} 开头的 net 在版图与 schematic 中都要有对应：

\begin{lstlisting}
same_nets!("TOP", "PAD*")
\end{lstlisting}

因此，若版图中有 PAD1 而 schematic 中没有，就是错误；若两边都有 PAD2 但彼此不匹配，同样是错误。

\ui{same\_nets} 与 \ui{same\_nets!} 可以出现在 LVS 脚本的任何位置。

关于 \ui{same\_nets} 的更多信息见 \ui{same\_nets} 与 \ui{same\_nets!}。

\subsection{circuit 等同提示}

默认情况下，同名的 circuit 被视为等价。若不成立，可以用 \ui{same\_circuit} 函数指定等价关系：

\begin{lstlisting}
same_circuits("CIRCUIT_IN_LAYOUT", "CIRCUIT_IN_SCHEMATIC")
\end{lstlisting}

把 circuit 声明为“相同”并不意味着跳过比较，它仍然会被比较；该函数只是提示到哪里去找比对目标。

\subsection{器件类等同提示}

默认情况下同名的器件类被视为等价。若不成立，可以用 \ui{same\_device\_classes} 指定：

\begin{lstlisting}
same_device_classes("PMOS_IN_LAYOUT", "PMOS_IN_SCHEMATIC") same_device_classes("NMOS_IN_LAYOUT", "NMOS_IN_SCHEMATIC")
\end{lstlisting}

该方法也可以多次调用，建立多对多的等价关系：

\begin{lstlisting}
same_device_classes("POLYRES", "RES") same_device_classes("WELLRES", "RES")
\end{lstlisting}

若某一侧给 \ui{nil}，则相应器件基本被忽略：

\begin{lstlisting}
# ignores "POLYRES" devices: same_device_classes("POLYRES", nil)
\end{lstlisting}

\subsection{容差}

比对器件参数时，默认要求严格相等。但对电阻这类器件，画出来的值通常很难与 schematic 中的值精确一致（尺寸会有取整、近似等误差）。

解决办法是指定参数容差。容差可以用绝对值或相对值给出：绝对容差表示版图参数可以在目标值上下偏差这么多；

相对容差则由目标值乘以该系数得出，例如 0.05 表示允许 5\% 的偏差。

也可以同时给出绝对与相对容差，此时两者相加，允许的偏差更大。

指定绝对容差用 \ui{tolerance} 函数：

\begin{lstlisting}
tolerance("NMOS", "L", 0.05)
\end{lstlisting}

两个参数分别是器件类名与要应用容差的参数名。上例中为 \ui{L} 指定了 50 纳米的容差。

相对容差作为第四个参数给出；可以把绝对容差设为 0，只保留相对容差。下面给出 1\% 的容差：

\begin{lstlisting}
tolerance("NMOS", "L", 0.0, 0.01)
\end{lstlisting}

容差还有更明确的写法：

\begin{lstlisting}
tolerance("NMOS", "L", :absolute => 0.05)
\end{lstlisting}

或者：

\begin{lstlisting}
tolerance("NMOS", "L", :relative => 0.01)
\end{lstlisting}

同时给出两者即表示绝对加相对容差。下面的调用为 \ui{NMOS} 的 \ui{L} 参数给出 50 纳米绝对容差与 1\% 相对容差。

\begin{lstlisting}
tolerance("NMOS", "L", 0.05, 0.01) tolerance("NMOS", "L", :absolute => 0.05, :relative => 0.01)
\end{lstlisting}

\subsection{忽略参数}

也可以让某些器件参数完全不参与 netlist 比对。例如不想比较 \ui{NMOS} 的 \ui{L} 参数，可以这样写：

\begin{lstlisting}
ignore_parameter("NMOS", "L")
\end{lstlisting}

这条语句可以放在脚本中 \ui{compare} 之前的任何位置。

默认只比较“主要”（primary）参数。以电阻为例，\ui{R} 是主要参数，\ui{L}、\ui{W}、\ui{A}、\ui{P} 都不是。用 \ui{tolerance} 启用某个参数会把它一并纳入比较。

\subsection{启用与禁用参数}

如前所述，某些器件参数是主要的，另一些不是。例如对电阻器件，\ui{R}（阻值）是主要参数，而 \ui{L}（长度）默认不是。

\begin{lstlisting}
enable_parameter("RES", "L")
\end{lstlisting}

这带来两点影响：其一，\ui{L} 参数会被写进 SPICE 输出 netlist；其二，它也会与 schematic 中的 \ui{L} 参数比较。

相应地，主要参数也可以用如下方式禁用：

\begin{lstlisting}
disable_parameter("RES", "R")
\end{lstlisting}

若该参数另有 \ui{tolerance} 或 \ui{ignore\_parameter} 设置，或安装了自定义的器件比较器，则上述行为会被覆盖。netlist 输出仅在该参数参与比较时才受影响。

\subsection{pin 互换}

当逻辑元件的输入在逻辑上等价、物理上却不同时，pin 互换（pin swapping）会很有用，例如 CMOS NAND 门：两个输入在逻辑上可交换，但版图上左右两个输入的实现并不相同。

\begin{lstlisting}
equivalent_pins("NAND_GATE", "A", "B")
\end{lstlisting}

第一个参数是版图 netlist 中的 circuit 名——只能在版图侧声明等价，不能对参考 schematic 声明。其后可以列出多个要互换的 pin。

\subsection{消去电容与电阻}

该功能可以消去“开路”的电阻与电容。目前串联电阻还不能被消去（短接）。

要消去所有阻值高于阈值的电阻，用 \ui{max\_res}。下面会消去所有阻值大于等于 1 kOhm 的电阻：

\begin{lstlisting}
max_res(1000)
\end{lstlisting}

要消去所有容值低于阈值的电容，用 \ui{min\_caps}。下面会消去所有容值小于等于 0.1 fF 的电容：

\begin{lstlisting}
min_caps(1e-16)
\end{lstlisting}

\subsection{检查 pin 标注}

LVS 本质上不依赖名字，除了消除歧义之外并不考虑 net 名：拓扑匹配优先。因此若 net 没有恰当地打 label，LVS 默认会把同拓扑但接错位置的 net 判为匹配。

这在涉及对外连接（例如 pad）时可能带来不良后果——芯片 pad 接错是致命错误。为缓解此问题，可用 \ui{flag\_mismatched\_nets} 功能检查名字不一致的 net。

该函数需要在比对步骤之后调用，即：

\begin{lstlisting}
compare flag_missing_ports
\end{lstlisting}

若只想在比较成功时抑制伪错误，也可以这样写：

\begin{lstlisting}
success = compare success && flag_missing_ports
\end{lstlisting}

该函数取 schematic 的顶层 circuit，检查所有连到 pin 的 net：核对对应版图 net 的名字（label），若名字不一致就报告。

在设计过程中自下而上推进时使用该功能，可以确保所有 pin 都被恰当标注。若你在顶层 circuit 中用 pin 来描述对外接口，这一点尤其重要。

注意，SPICE netlist 的顶层 circuit 也可能没有 pin（例如顶层不是 SUBCKT 时），这种情况下该函数不会做检查。

你可以使用

\begin{lstlisting}
schematic.make_top_level_pins
\end{lstlisting}

在没有 pin 的情况下创建 pin。但这样会把每条 net 都变成 pin，从而强迫你标注顶层 circuit 中的每一条 net，因此更好的做法还是显式提供 pin。

\subsection{比对与网表层次}

好的版图是按层次构建的，netlist 比对也能利用层次结构。所谓“按层次”，是指 circuit 由若干 Cell 构成，而这些 Cell 又对应 subcircuit，形成与版图一致的多级结构。

给定层次化的版图与 schematic netlist，比对算法将自下而上进行：先比较叶 circuit（不含 subcircuit 调用的 circuit），若它们匹配，再继续向上。

要启用层次化提取，必须使用 \ui{deep} 模式（\ui{deep}）；若缺少该语句，版图 netlist 就是平坦的（不含 subcircuit）。

（\ui{align}）：它会把在另一侧 netlist 中不存在的 circuit 从版图或 schematic netlist 中移除。版图中常有只起辅助作用的 Cell，它们在 schematic 中并没有对应电路。

一般建议把 \ui{align} 放在 \ui{netlist.simplify} 之类 netlist 处理以及 \ui{compare} 之前。

（\ui{align}）还有一个很有用的副作用：它会去掉任一侧顶层 circuit 之上的 circuit，从而最终从 circuit 树中切出一棵子树，只保留两侧共有的部分。

\subsection{比对算法的工作方式}

netlist 比对算法的大致流程如下：

\begin{lstlisting}
foreach circuit bottom up: if matching circuit found in reference netlist: if all subcircuits have been matched and pin matching has been established for them: compare net graph locally from this circuit else: skip circuit with warning else: issue a circuit mismatch error
\end{lstlisting}

这一流程的后果是：一旦某个 circuit 的 pin 无法与参考 circuit 的 pin 对应，或无法与该层次上的其它 circuit 对应，比对就不再处理其父级 circuit。

局部 net 图比对算法是一种回溯算法，并借助拓扑 net 分类进行提示。拓扑分类基于最近 net 邻域，也就是从某条 net 出发，经过器件与连线能到达哪些 net、各需几跳。

\klfig[0.8\textwidth]{net_graph.png}{}

这里 IN net 的邻域是：经“栅极到源漏、越过 M1”的一跳到达 VDD；经“栅极到源漏、越过 M1 与 M2”的两跳到达 OUT；再经另一条路径到达 VSS。

在大型 netlist 中，算法会先按邻域度量尝试无歧义地匹配 net 并登记为成对 net，这些配对往往足以推断出其余 net 的对应关系，从而大幅缩小搜索空间。

\section{LVS 网表调整}\label{sec:lvs_tweaks}

netlist 调整（tweaking）对于把 netlist 标准化很重要。若不做调整，提取出的 netlist 里可能含有冗余元素，或含有 schematic 中根本找不到对应的东西，比对就很难通过。

netlist 调整作用于已提取的 netlist 对象，可用 \ui{netlist} 函数取得；该函数若尚未提取会先触发提取。

调整也可以作用在 schematic netlist 上，例如要把名为 \ui{NMOS} 的模型 subcircuit 打平，可以这样写：

\begin{lstlisting}
schematic.flatten_circuit("NMOS")
\end{lstlisting}

\subsection{生成顶层 pin}

提取出的 circuit 在顶层没有 pin，因为提取器无从得知该 circuit 对外接到哪里。比对函数不会去猜，因此顶层若无 pin，就只能做拓扑层面的比对。

KLayout 提供一个用简单启发式建立顶层 pin 的函数：顶层 circuit 中每一条有名字（有 label）的 net 都会生成一个 pin：

\begin{lstlisting}
netlist.make_top_level_pins
\end{lstlisting}

\subsection{器件合并}

器件合并对于那些在版图中并非以单一实体表示的器件很重要，例如：

\begin{itemize}
\item 指状 MOS 晶体管：宽度很大的 MOS 晶体管常被拆成多段，以降低栅极与扩散区的寄生电阻、电容。
\item 串联电阻：大电阻常被分成多条、再以折线（meander）方式串起来。从器件角度看，这样的电阻由若干段串联电阻构成。
\item 阵列电容：大电容常被拆成多个小电容排成阵列并联，以控制寄生串联电阻。
\end{itemize}

这些情况下，schematic 通常把它们归并成单个器件，用集中参数值（总电阻、总电容、晶体管宽度）表示，因此提取网表也要做同样的合并。

\begin{lstlisting}
netlist.combine_devices
\end{lstlisting}

合并串联器件可能留下悬空 net（原本连接这些器件的 net）。这类 net 可以用 \ui{purge} 移除。见 \ui{purge}，它是若干方法的封装。

建议在这一步之前先运行 \ui{make\_toplevel\_pins} 与 \ui{purge}（见下文）。

\subsection{打平 circuit（消去层次）}

经常需要打平那些并不代表某个组织层次的 circuit——它们只是别的什么东西的外壳。版图中器件常常由 PCell 或单元实现，这些在 schematic 中并没有对应层次。

要从提取网表中打平某个 circuit，用：

\begin{lstlisting}
netlist.flatten_circuit("CIRCUIT_NAME")
\end{lstlisting}

\ui{flatten\_circuit} 的参数是 glob 通配符（类 shell）。例如 \ui{NMOS*} 会打平所有以 \ui{NMOS} 开头的 circuit。

\subsection{自动打平 circuit（网表对齐）}

除了显式打平，也可以通过 \ui{align} 方法自动打平。

这一步是可选的，但很有用：它会找出没有对应 schematic circuit 的版图 Cell 并打平它们。“打平”意味着其内容被并入父级，本身不再作为一层出现。

\ui{align} 同样会打平那些没有对应版图 Cell 的 schematic circuit。

\begin{lstlisting}
align
\end{lstlisting}

\subsection{黑盒化（circuit 抽象）}

circuit 抽象是一种降低验证开销的手段。在早期阶段，把某个 Cell 替换成简化版本或只留下 pin 的框架往往很有用。

与此相关的一个有用方法是 \ui{blank\_circuit}：清空某 circuit 的内部结构。此后比对算法会把两边都视为黑盒，只比较接口。

要清空某个 circuit 的内部，用该方法：

\begin{lstlisting}
netlist.blank_circuit("CIRCUIT_NAME") schematic.blank_circuit("CIRCUIT_NAME")
\end{lstlisting}

注意：在本版本中，blank\_circuit 要在 purge 或 simplify 之前使用（均为 LVS 脚本函数）。blank\_circuit 会设置一个标志，阻止抽象 circuit 被清除。

它也有简写形式（\ui{blank\_circuit}）。与基于 netlist 的 \ui{blank\_circuit} 不同，这个简写可以用在 LVS 脚本的任何位置：

\begin{lstlisting}
blank_circuit("CIRCUIT_NAME")
\end{lstlisting}

两种情况下 \ui{blank\_circuit} 的参数都是 glob 通配符（类 shell）。例如 \ui{MEMORY*} 会清空所有以 \ui{MEMORY} 开头的 circuit。

\subsection{对称节点合并}

有时版图里可以省掉一些连接，因为它们不承载电流。这能简化版图、允许更紧凑的排布，但最终会与 schematic 的连通性不一致。

来解决 \ui{split\_gates} 问题（见下文 \ui{split\_gates}）：在所示配置中，内部源漏节点是对称的，因此可以让 LVS 自动把它们视为同一节点。

\begin{lstlisting}
join_symmetric_nets("NAND2")
\end{lstlisting}

不过还有一个更专门的特性（\ui{split\_gates}），它覆盖更多情形，但专门针对 MOS 器件。

\subsection{分裂栅}

下图描述了这种称为“分裂栅”（split gate）的情形。此例中 N1 与 N2 完全相同，互换它们不会改变电路的功能。

\klfig[0.92\textwidth]{lvs_symm_nodes.png}{}

KLayout 提供一个特性（\ui{split\_gates}），会在 netlist 提取之后补上这类连接：

\begin{lstlisting}
split_gates("NMOS")
\end{lstlisting}

该函数会分析提取网表中的所有 circuit，针对 \ui{NMOS} 器件，把内部相关的分裂栅源漏节点连起来。若该函数找不到可连的分裂栅，则不影响结果。

器件名必须是 MOS3、MOS4、DMOS3 或 DMOS4 器件。所有相关器件的栅长必须完全一致；对 MOS4 与 DMOS4，同一栅极 net 上的所有器件还必须共享同一个 bulk net。

除器件名之外，还可以给一个 glob 风格的 circuit 通配符，此时分析只限于匹配该模式的 circuit。

\ui{split\_gates} 可以用在 LVS 脚本的任何位置。

\subsection{清除冗余}

提取网表中常含没有功能意义的元素：例如 via 单元会生成只有一个 pin、没有任何器件的 subcircuit；孤立的金属岛也会形成不含器件的 subcircuit。

（\ui{purge}）会移除所有悬空 net 与所有不含器件或 subcircuit 的 circuit；\ui{purge\_nets} 只清除悬空 net。悬空 net 指不连接任何器件或 subcircuit 的 net。

\begin{lstlisting}
netlist.purge netlist.purge_nets
\end{lstlisting}

\subsection{归一化封装（simplify）}

\ui{simplify} 是 \ui{make\_top\_level\_pins}、\ui{purge}、\ui{combine\_devices} 与 \ui{purge\_nets} 的封装，按如下推荐顺序执行：

\begin{lstlisting}
netlist.simplify
\end{lstlisting}

技术细节上，之所以把这序列里包含 \ui{make\_top\_level\_pins}，是因为有 pin 之后 net 就不算悬空，这样 \ui{purge\_nets} 会保留有 label 的 net 的 pin。
